\chapter{Conclusions}

\section{Summary of Contributions}

This thesis addressed a concrete gap in the RISC-V custom-extension ecosystem: at
the outset, the Composable eXtension Unit (CXU) specification existed as a design
document but had no hardware implementation on any processor, and no sandbox
environment supported its features. This work delivers both.

The \textbf{CXU Playground} is an open-source, full-stack development and
benchmarking environment for RISC-V custom extensions, built on the LiteX SoC
framework and targeting the VexiiRiscv processor. The following components were
designed, implemented, and validated as part of this work:

\begin{itemize}
    \item A \textbf{SpinalHDL pipeline plugin} for VexiiRiscv that implements the
          CXU interface: opcode interception at the execute stage, forwarding of
          \texttt{rs1}/\texttt{rs2} operands, writeback of the result to \texttt{rd},
          pipeline stall management for multi-cycle operations, and
          \texttt{cxidx}/\texttt{cxdata} CSR support for extension selection and
          shared state.

    \item A \textbf{Buildroot-integrated Linux environment} that cross-compiles the
          Linux kernel, root filesystem, and user-space utilities for the RV32IMAFC
          target, enabling OS-level benchmarking under realistic system conditions.

    \item A \textbf{C-level CXU runtime library} that wraps the assembly-level
          invocation into safe callable functions, manages extension multiplexing
          via the \texttt{cxidx} register, and provides a portable API that
          eliminates the need for inline assembly in user code.

    \item Three \textbf{case-study extensions} --- AES GF$(2^8)$ acceleration,
          GZIP/DEFLATE bit-manipulation acceleration, and a TFLite Micro
          Multiply-Accumulate (MAC) extension --- each with a paired software
          reference implementation, a Verilator verification testbench, and a
          benchmarking suite.
\end{itemize}

\section{Key Results}

\textbf{AES acceleration} produced the most dramatic result. Offloading SubBytes
and MixColumns arithmetic to a single-cycle combinational CXU instruction reduced
cycles per encrypted byte from approximately 2{,}670 to approximately 506, a
consistent \textbf{5.25$\times$ speedup} across all tested input sizes from 128 B
to 16{,}768 B. The gain is directly attributable to replacing 5--6 iteration loops
of 3--4 instructions each with a single pipeline cycle per $\text{GF}(2^8)$
operation. The speedup is stable across input sizes, confirming that fixed
invocation overhead is negligible relative to the per-byte compute savings. OS
overhead in the Linux environment adds only $\pm$2\% jitter with no measurable
change in the mean speedup factor.

\textbf{GZIP decompression} showed an approximately \textbf{10\% reduction in CPU
time} for the bit-stream manipulation component of DEFLATE decompression. The
smaller end-to-end gain compared to AES reflects two factors: the LZ77 sliding
window search --- a substantial fraction of total decompression work --- is not
accelerated by the current extension, and the end-to-end measurements include
filesystem I/O, which dilutes the compute-side improvement. Profiling confirms that
this result is consistent with the fraction of the inner loop covered by the
four custom instructions.

\textbf{TFLite Micro MAC} extension correctness was verified against the TFLite
Micro built-in inference self-tests; hardware timing data collection is ongoing.

Together, these results confirm that the CXU Playground can host extensions of
varying complexity and that even a small number of tightly coupled custom
instructions can yield order-of-magnitude cycle reductions for
arithmetic-intensive workloads.

\section{Limitations of the Current Implementation}

Three limitations of the current implementation are worth highlighting.

First, \textbf{CSR access latency} is higher than the CXU specification assumes.
Because \texttt{cxidx} and \texttt{cxdata} are memory-mapped over the Wishbone
bus in the LiteX SoC rather than being native pipeline registers, each state read
or write incurs a full bus-transaction penalty. This makes the shared-state
mechanism best suited to one-time configuration rather than high-bandwidth
inter-invocation communication, and penalises workloads that switch frequently
between multiple extensions.

Second, the \textbf{Linux kernel does not yet save and restore CXU state} as part
of the task context switch. In a preemptive, multi-process workload, this is a
correctness issue: if two processes share a physical CXU and the scheduler
preempts one mid-invocation, the extension state can be corrupted. The playground
is therefore currently safe only for single-process CXU workloads under Linux.

Third, \textbf{CXU capabilities are not described in the Linux Device Tree}: the
runtime library hard-codes extension addresses rather than discovering them
dynamically through the kernel's device model.

\section{Future Work}

The CXU Playground provides a solid foundation for several directions of future
work.

\textbf{Kernel context-switch support} is the most pressing correctness gap. The
CXU specification defines the required save/restore semantics; implementing them
requires adding \texttt{cxidx}/\texttt{cxdata} to the task struct and instrumenting
the scheduler's context-switch path. This work would make the playground safe for
multi-process workloads and would complete the OS-facing portion of the CXU
specification.

\textbf{Device Tree integration} should follow naturally: extending the LiteX SoC
generator to emit DTS nodes describing available extensions (index range, state
width, memory-access permissions) would allow a kernel driver to enumerate and
manage extensions through the standard Linux device model, and would enable
future hardware/software co-design tools to query the platform capabilities
programmatically.

\textbf{CSR latency reduction} could be addressed by implementing \texttt{cxidx}
and \texttt{cxdata} as native VexiiRiscv pipeline registers rather than
Wishbone-mapped registers, bringing access latency in line with the CXU
specification's assumptions and enabling stateful extensions to operate efficiently.

\textbf{Post-quantum cryptography} is a natural extension of the case studies in
Chapter~5. The related work survey in Chapter~2 documents speedups of up to
9.6$\times$ for CRYSTALS-Kyber using tightly coupled custom instructions
\cite{Levi2020RISQV}. Implementing NTT butterfly and modular reduction operations
as CXU extensions, and measuring them within the Playground, would both validate
the framework for a security-critical workload and produce results directly
comparable to RISQ-V and similar prior work.

\textbf{Cross-core portability} --- porting the CXU plugin to a second RV32
implementation such as CVA6 or the SERV core --- would validate the original
premise of the composable-extensions family: that a single extension implementation
can be reused across different processor microarchitectures. This would also test
whether the SpinalHDL plugin abstraction is genuinely separable from VexiiRiscv's
internals.

\textbf{FPGA resource and Fmax characterisation} should be added to the benchmarking
workflow. Reporting LUT count, flip-flop usage, BRAM consumption, and synthesis
Fmax alongside cycle-count data would give extension designers an accurate picture
of the area-performance trade-off and would make the Playground's benchmark results
directly comparable to FPGA-based results in the related-work literature.

\section{Closing Remarks}

The CXU Playground demonstrates that a full-stack, Linux-capable sandbox for
composable RISC-V custom extensions is practically achievable on commodity FPGA
hardware. The framework successfully lowers the barrier to entry for custom
hardware acceleration research by automating SoC generation, Linux integration,
cross-compilation, and benchmarking into a single cohesive environment. The
measured 5.25$\times$ AES speedup confirms that even a small number of carefully
designed, tightly coupled hardware instructions can deliver substantial gains for
the class of workloads most relevant to embedded cryptography and machine learning.
The open-source playground provides a reproducible platform for the community to
build on, and the limitations identified here chart a clear path toward a more
complete implementation.
